\documentclass[11pt]{article}
% Shared preamble for the printable mini lessons (1.4 onward).
%   \documentclass[11pt]{article}
%   % Shared preamble for the printable mini lessons (1.4 onward).
%   \documentclass[11pt]{article}
%   % Shared preamble for the printable mini lessons (1.4 onward).
%   \documentclass[11pt]{article}
%   \input{../preamble}
% Build the key with:  pdflatex -jobname=<name>_KEY "\def\iskey{}\input{<name>}"
% Figures shared with the web page live in figures/Unit_X/<name>.tex and are
% pulled in with \fig{Unit_X/<name>}; the standalone package drops their own
% preamble, so their macros come from figures/preamble.tex, loaded here.
\input{../../figures/preamble}
\usepackage[margin=0.75in]{geometry}
\usepackage{array}
\usepackage{enumitem}
\usepackage{amsmath}
\usepackage[helvet]{sfmath}  % math in the same sans font as the text
\usepackage{standalone}
\pagestyle{empty}
\setlength{\parindent}{0pt}
\setlist[enumerate,1]{leftmargin=1.5em,itemsep=1em}
\setlist[enumerate,2]{label=(\alph*),leftmargin=2em,itemsep=2.2em}

\newif\ifkey
\ifdefined\iskey \keytrue \fi
\colorlet{keyred}{red!75!black}
\newcommand{\ans}[1]{\ifkey\leavevmode{\color{keyred}#1}\fi}
% vertical work space on the student copy, less on the key
\newcommand{\work}[1]{\ifkey\vspace{0.6em}\else\vspace{#1}\fi}
\newcommand{\fig}[1]{\input{../../figures/#1}}
% Title line: \lessontitle{Quartiles, IQR, and Box Plots}
\newcommand{\lessontitle}[1]{{\Large\bfseries #1\ifkey\ \textcolor{keyred}{KEY}\fi}\hfill Name \underline{\hspace{2in}}\par\vspace{0.6em}}
\newcommand{\partner}{\par\vspace{0.4em}Work with a partner. Compare answers when you finish each problem.\par}

% Build the key with:  pdflatex -jobname=<name>_KEY "\def\iskey{}\input{<name>}"
% Figures shared with the web page live in figures/Unit_X/<name>.tex and are
% pulled in with \fig{Unit_X/<name>}; the standalone package drops their own
% preamble, so their macros come from figures/preamble.tex, loaded here.
% Shared preamble for web figures.  Each figure is a standalone .tex that
% \input's this file; scripts/make-figures.sh turns it into an SVG.
\usepackage[T1]{fontenc}
\usepackage{helvet}
\renewcommand{\familydefault}{\sfdefault}
\usepackage{tikz}
\usetikzlibrary{decorations.pathreplacing}
% axis numbers in the text font, not math italic
\pgfkeys{/pgf/number format/.cd,assume math mode=true,1000 sep={}}

% dot plot: \dotplot{min}{max}{value/count,...}{axis label}
\newcommand{\dotplot}[4]{%
\begin{tikzpicture}[x=0.9cm,y=0.33cm]
  \draw[thick] (#1-0.5,0) -- (#2+0.5,0);
  \foreach \t in {#1,...,#2} {\draw (\t,0.25) -- (\t,-0.25) node[below] {\t};}
  \foreach \v/\c in {#3} {\ifnum\c>0 \foreach \k in {1,...,\c} {\fill (\v,\k*1.15) circle (0.16cm);}\fi}
  \node at ({(#1+#2)/2},-2.3) {#4};
\end{tikzpicture}}

% shape sketch histogram (12 bins): \shapehist{heights}{label}
\newcommand{\shapehist}[2]{%
  \foreach \h [count=\i from 0] in {#1} \draw[fill=gray!30] (\i,0) rectangle (\i+1,\h);
  \draw (-0.3,0) -- (12.3,0);
  \node[font=\bfseries] at (6,-2.6) {#2};}
\newcommand{\meanmark}[1]{\fill (#1,0) -- ++(-0.3,-0.75) -- ++(0.6,0) -- cycle;}
\newcommand{\medmark}[2]{\draw[very thick,dashed] (#1,0) -- (#1,#2);}

% number line with minor ticks: \nline{min}{max}{label step}{axis label}
\newcommand{\nline}[4]{%
  \draw[thick] (#1-0.5,0) -- (#2+0.5,0);
  \foreach \t in {#1,...,#2} \draw (\t,0.12) -- (\t,-0.12);
  \pgfmathtruncatemacro{\nxt}{#1+#3}
  \foreach \t in {#1,\nxt,...,#2} \draw (\t,0.2) -- (\t,-0.2) node[below] {\t};
  \node at ({(#1+#2)/2},-0.95) {#4};}
% box plot above the line: \bplot{min}{Q1}{med}{Q3}{max}
\newcommand{\bplot}[5]{%
  \draw[thick] (#1,1.0) -- (#2,1.0);
  \draw[thick] (#4,1.0) -- (#5,1.0);
  \draw[thick] (#1,0.75) -- (#1,1.25);
  \draw[thick] (#5,0.75) -- (#5,1.25);
  \draw[thick,fill=gray!20] (#2,0.55) rectangle (#4,1.45);
  \draw[very thick] (#3,0.55) -- (#3,1.45);}

% ---- added for 1.4 onward ----
% axis on y=0: \axis{min}{max}{label step}{tick step}{axis label}
\newcommand{\axis}[5]{%
  \draw[thick] (#1,0) -- (#2,0);
  \pgfmathsetmacro{\axT}{#1+#4}\foreach \t in {#1,\axT,...,#2} \draw (\t,0.1) -- (\t,-0.1);
  \pgfmathsetmacro{\axL}{#1+#3}\foreach \t in {#1,\axL,...,#2}
    \draw (\t,0.18) -- (\t,-0.18) node[below] {\pgfmathprintnumber{\t}};
  \node at ({(#1+#2)/2},-0.95) {#5};}
% box plot centred at height y: \bplotat{y}{min}{Q1}{med}{Q3}{max}
\newcommand{\bplotat}[6]{%
  \draw[thick] (#2,#1) -- (#3,#1);
  \draw[thick] (#5,#1) -- (#6,#1);
  \draw[thick] (#2,#1-0.22) -- (#2,#1+0.22);
  \draw[thick] (#6,#1-0.22) -- (#6,#1+0.22);
  \draw[thick,fill=gray!20] (#3,#1-0.42) rectangle (#5,#1+0.42);
  \draw[very thick] (#4,#1-0.42) -- (#4,#1+0.42);}
% an outlier dot: \outlier{x}{y}
\newcommand{\outlier}[2]{\fill (#1,#2) circle (2.4pt);}

\usepackage[margin=0.75in]{geometry}
\usepackage{array}
\usepackage{enumitem}
\usepackage{amsmath}
\usepackage[helvet]{sfmath}  % math in the same sans font as the text
\usepackage{standalone}
\pagestyle{empty}
\setlength{\parindent}{0pt}
\setlist[enumerate,1]{leftmargin=1.5em,itemsep=1em}
\setlist[enumerate,2]{label=(\alph*),leftmargin=2em,itemsep=2.2em}

\newif\ifkey
\ifdefined\iskey \keytrue \fi
\colorlet{keyred}{red!75!black}
\newcommand{\ans}[1]{\ifkey\leavevmode{\color{keyred}#1}\fi}
% vertical work space on the student copy, less on the key
\newcommand{\work}[1]{\ifkey\vspace{0.6em}\else\vspace{#1}\fi}
\newcommand{\fig}[1]{\input{../../figures/#1}}
% Title line: \lessontitle{Quartiles, IQR, and Box Plots}
\newcommand{\lessontitle}[1]{{\Large\bfseries #1\ifkey\ \textcolor{keyred}{KEY}\fi}\hfill Name \underline{\hspace{2in}}\par\vspace{0.6em}}
\newcommand{\partner}{\par\vspace{0.4em}Work with a partner. Compare answers when you finish each problem.\par}

% Build the key with:  pdflatex -jobname=<name>_KEY "\def\iskey{}\input{<name>}"
% Figures shared with the web page live in figures/Unit_X/<name>.tex and are
% pulled in with \fig{Unit_X/<name>}; the standalone package drops their own
% preamble, so their macros come from figures/preamble.tex, loaded here.
% Shared preamble for web figures.  Each figure is a standalone .tex that
% \input's this file; scripts/make-figures.sh turns it into an SVG.
\usepackage[T1]{fontenc}
\usepackage{helvet}
\renewcommand{\familydefault}{\sfdefault}
\usepackage{tikz}
\usetikzlibrary{decorations.pathreplacing}
% axis numbers in the text font, not math italic
\pgfkeys{/pgf/number format/.cd,assume math mode=true,1000 sep={}}

% dot plot: \dotplot{min}{max}{value/count,...}{axis label}
\newcommand{\dotplot}[4]{%
\begin{tikzpicture}[x=0.9cm,y=0.33cm]
  \draw[thick] (#1-0.5,0) -- (#2+0.5,0);
  \foreach \t in {#1,...,#2} {\draw (\t,0.25) -- (\t,-0.25) node[below] {\t};}
  \foreach \v/\c in {#3} {\ifnum\c>0 \foreach \k in {1,...,\c} {\fill (\v,\k*1.15) circle (0.16cm);}\fi}
  \node at ({(#1+#2)/2},-2.3) {#4};
\end{tikzpicture}}

% shape sketch histogram (12 bins): \shapehist{heights}{label}
\newcommand{\shapehist}[2]{%
  \foreach \h [count=\i from 0] in {#1} \draw[fill=gray!30] (\i,0) rectangle (\i+1,\h);
  \draw (-0.3,0) -- (12.3,0);
  \node[font=\bfseries] at (6,-2.6) {#2};}
\newcommand{\meanmark}[1]{\fill (#1,0) -- ++(-0.3,-0.75) -- ++(0.6,0) -- cycle;}
\newcommand{\medmark}[2]{\draw[very thick,dashed] (#1,0) -- (#1,#2);}

% number line with minor ticks: \nline{min}{max}{label step}{axis label}
\newcommand{\nline}[4]{%
  \draw[thick] (#1-0.5,0) -- (#2+0.5,0);
  \foreach \t in {#1,...,#2} \draw (\t,0.12) -- (\t,-0.12);
  \pgfmathtruncatemacro{\nxt}{#1+#3}
  \foreach \t in {#1,\nxt,...,#2} \draw (\t,0.2) -- (\t,-0.2) node[below] {\t};
  \node at ({(#1+#2)/2},-0.95) {#4};}
% box plot above the line: \bplot{min}{Q1}{med}{Q3}{max}
\newcommand{\bplot}[5]{%
  \draw[thick] (#1,1.0) -- (#2,1.0);
  \draw[thick] (#4,1.0) -- (#5,1.0);
  \draw[thick] (#1,0.75) -- (#1,1.25);
  \draw[thick] (#5,0.75) -- (#5,1.25);
  \draw[thick,fill=gray!20] (#2,0.55) rectangle (#4,1.45);
  \draw[very thick] (#3,0.55) -- (#3,1.45);}

% ---- added for 1.4 onward ----
% axis on y=0: \axis{min}{max}{label step}{tick step}{axis label}
\newcommand{\axis}[5]{%
  \draw[thick] (#1,0) -- (#2,0);
  \pgfmathsetmacro{\axT}{#1+#4}\foreach \t in {#1,\axT,...,#2} \draw (\t,0.1) -- (\t,-0.1);
  \pgfmathsetmacro{\axL}{#1+#3}\foreach \t in {#1,\axL,...,#2}
    \draw (\t,0.18) -- (\t,-0.18) node[below] {\pgfmathprintnumber{\t}};
  \node at ({(#1+#2)/2},-0.95) {#5};}
% box plot centred at height y: \bplotat{y}{min}{Q1}{med}{Q3}{max}
\newcommand{\bplotat}[6]{%
  \draw[thick] (#2,#1) -- (#3,#1);
  \draw[thick] (#5,#1) -- (#6,#1);
  \draw[thick] (#2,#1-0.22) -- (#2,#1+0.22);
  \draw[thick] (#6,#1-0.22) -- (#6,#1+0.22);
  \draw[thick,fill=gray!20] (#3,#1-0.42) rectangle (#5,#1+0.42);
  \draw[very thick] (#4,#1-0.42) -- (#4,#1+0.42);}
% an outlier dot: \outlier{x}{y}
\newcommand{\outlier}[2]{\fill (#1,#2) circle (2.4pt);}

\usepackage[margin=0.75in]{geometry}
\usepackage{array}
\usepackage{enumitem}
\usepackage{amsmath}
\usepackage[helvet]{sfmath}  % math in the same sans font as the text
\usepackage{standalone}
\pagestyle{empty}
\setlength{\parindent}{0pt}
\setlist[enumerate,1]{leftmargin=1.5em,itemsep=1em}
\setlist[enumerate,2]{label=(\alph*),leftmargin=2em,itemsep=2.2em}

\newif\ifkey
\ifdefined\iskey \keytrue \fi
\colorlet{keyred}{red!75!black}
\newcommand{\ans}[1]{\ifkey\leavevmode{\color{keyred}#1}\fi}
% vertical work space on the student copy, less on the key
\newcommand{\work}[1]{\ifkey\vspace{0.6em}\else\vspace{#1}\fi}
\newcommand{\fig}[1]{\input{../../figures/#1}}
% Title line: \lessontitle{Quartiles, IQR, and Box Plots}
\newcommand{\lessontitle}[1]{{\Large\bfseries #1\ifkey\ \textcolor{keyred}{KEY}\fi}\hfill Name \underline{\hspace{2in}}\par\vspace{0.6em}}
\newcommand{\partner}{\par\vspace{0.4em}Work with a partner. Compare answers when you finish each problem.\par}

% Unit summary: the unit's tools in tables, to print and keep. Same tables as
% chapters/Unit_2/Summary.qmd; edit both together. Nothing to answer, so the
% _KEY copy is the same sheet.
\newcommand{\sumhead}[1]{\par\vspace{0.35em}{\large\bfseries #1}\par\vspace{0.15em}}
\newcolumntype{L}[1]{>{\raggedright\arraybackslash}p{#1}}
\renewcommand{\arraystretch}{1.1}
% room above and below a full-size fraction in a table row
\newcommand{\tall}{\rule[-0.85em]{0pt}{2.2em}}

\begin{document}

\lessontitle{Unit 2 Summary: Analyzing One Variable}

\sumhead{Standard Deviation}
\begin{tabular}{|L{1.45in}|L{5.2in}|}\hline
\textbf{Standard deviation} & The SD is about how far the values usually are from the mean.\\ \hline
\textbf{Small SD} & The values are close to the mean. Player A scores close to 10 every game.\\ \hline
\textbf{Large SD} & The values are spread out. Player B might score 2 or 18.\\ \hline
\textbf{Facts} & The SD is never negative. It is 0 only when every value is the same. It has the same units as the data (points, minutes, inches).\\ \hline
\end{tabular}

\vspace{0.5em}
\begin{tabular}{|L{2.75in}|L{3.9in}|}\hline
\textbf{Step} & \textbf{Example: 7, 7, 10, 13, 13}\\ \hline
1. Find the mean, $\bar{x}$. & $\bar{x} = 50 \div 5 = 10$\\ \hline
2. Subtract the mean from each value. & $-3, -3, 0, 3, 3$\\ \hline
3. Square each result. & $9, 9, 0, 9, 9$\\ \hline
4. Add the squares. & $\Sigma(x - \bar{x})^2 = 9 + 9 + 0 + 9 + 9 = 36$\\ \hline
5. Divide by $n - 1$. This is the \textbf{variance}. & $36 \div 4 = 9$\\ \hline
6. Take the square root. & $s = \sqrt{9} = 3$\\ \hline
\tall\textbf{The formula:} \; $s = \sqrt{\dfrac{\Sigma(x - \bar{x})^2}{n - 1}}$ &
$\Sigma$ (capital sigma) means \textbf{add up}. $\Sigma(x - \bar{x})^2$ means: find $(x - \bar{x})^2$ for every value, then add them all. $n$ is how many values there are. $s$ is the standard deviation.\\ \hline
\end{tabular}

\sumhead{Choosing Measures}
\begin{tabular}{|L{2.45in}|L{2.0in}|L{2.0in}|}\hline
\textbf{Shape of the data} & \textbf{Center} & \textbf{Spread}\\ \hline
Symmetric, no outliers & mean & standard deviation\\ \hline
Skewed, or has outliers & median & IQR\\ \hline
\end{tabular}

\vspace{0.3em}
\textbf{Why:} an outlier pulls the mean and the SD. It does not pull the median and the IQR much.

\sumhead{Comparing Two Distributions}
\begin{tabular}{|L{0.85in}|L{2.05in}|L{3.55in}|}\hline
 & \textbf{Ask} & \textbf{Example: hours of sleep last night}\\ \hline
\textbf{Shape} & Symmetric or skewed? & 9th grade is symmetric. 12th grade is skewed right.\\ \hline
\textbf{Center} & Which group is higher? (median or mean) & The median for 9th grade (8 hours) is greater than for 12th grade (6 hours). 9th graders usually sleep about 2 hours more.\\ \hline
\textbf{Spread} & Which group is more spread out? (IQR or SD) & Both groups have an IQR of 2 hours. The spreads are about the same.\\ \hline
\textbf{Outliers} & Any values far from the rest? & One 12th grader slept 12 hours. That is an outlier (the upper fence is 10).\\ \hline
\end{tabular}

\vspace{0.3em}
\textbf{Use comparing words:} ``The median for 9th grade is \textbf{2 hours greater than} for 12th grade.'' Use numbers and the context.

\sumhead{Percentiles}
\begin{tabular}{|L{3.25in}|L{3.4in}|}\hline
\textbf{Percentile} & \textbf{Example: 20 students took a test}\\ \hline
The percent of the data that is \textbf{below} the value.\newline
\tall$\dfrac{\text{number of values below}}{\text{total number of values}} \times 100$ &
15 students scored below 85. $15 \div 20 = 0.75$. A score of 85 is at the \textbf{75th percentile}. It is higher than 75\% of the class.\\ \hline
A percentile is \textbf{not} a percent score. &
Scoring at the 75th percentile does \textbf{not} mean you got 75\% of the questions right. It means you scored higher than 75\% of the people.\\ \hline
\end{tabular}

\newpage
\sumhead{Z-Scores}
\begin{tabular}{|L{1.2in}|L{2.75in}|L{2.5in}|}\hline
 & \textbf{What it tells} & \textbf{Example: mean 70, SD 10}\\ \hline
\textbf{Z-score}\newline \tall$z = \dfrac{x - \mu}{\sigma}$ &
How many SDs a value is above or below the mean. $\mu$ (mu) is the mean. $\sigma$ (small sigma) is the SD. &
\tall$z = \dfrac{85 - 70}{10} = 1.5$\newline 85 is \textbf{1.5 SDs above} the mean.\\ \hline
\textbf{Sign of $z$} & $z > 0$: above the mean. \newline $z < 0$: below the mean. \; $z = 0$: equal to the mean. &
55 has $z = -1.5$. It is 1.5 SDs \textbf{below} the mean.\\ \hline
\textbf{Going backward} & $x = \mu + z\sigma$ & $z = -2$: \; $70 + (-2)(10) = 50$\\ \hline
\textbf{Comparing} & Z-scores let you compare values from \textbf{different} data sets. &
Ana: 85, $z = 1.5$. Ben: 80 (mean 74, SD 3), $z = 2$. Ben did better compared to his class.\\ \hline
\end{tabular}

\sumhead{Normal Distributions and the Empirical Rule}
\begin{minipage}[c]{3.3in}\centering
\scalebox{0.6}{\fig{Unit_2/2.4_empirical}}
\end{minipage}\hfill
\begin{minipage}[c]{3.45in}\raggedright
\begin{tabular}{|L{0.6in}|L{1.5in}|L{0.6in}|}\hline
\textbf{Within} & \textbf{From} & \textbf{About}\\ \hline
1 SD & $\mu - \sigma$ to $\mu + \sigma$ & 68\%\\ \hline
2 SDs & $\mu - 2\sigma$ to $\mu + 2\sigma$ & 95\%\\ \hline
3 SDs & $\mu - 3\sigma$ to $\mu + 3\sigma$ & 99.7\%\\ \hline
\end{tabular}

\vspace{0.3em}
A \textbf{normal distribution} is bell-shaped and symmetric. The mean and the median are in the center (they are equal). Two numbers describe it: the mean $\mu$ (center) and the SD $\sigma$ (spread).
\end{minipage}

\vspace{0.3em}
\begin{tabular}{|L{3.75in}|L{2.9in}|}\hline
\textbf{Example: heights, mean 64 inches, SD 3 inches} & \textbf{Answer}\\ \hline
Percent from 61 to 67 inches & 68\% (within 1 SD)\\ \hline
Percent from 58 to 70 inches & 95\% (within 2 SDs)\\ \hline
Percent taller than 67 inches & $13.5\% + 2.35\% + 0.15\% = 16\%$\\ \hline
Percentile of 67 inches & $100\% - 16\% = 84\%$, the 84th percentile\\ \hline
\end{tabular}

\vspace{0.3em}
\textbf{Changes:} a larger SD makes the curve wider and flatter. A different mean moves the curve left or right.

\sumhead{Normal Calculations with a Z-Table}
\textbf{The z-table gives the area to the left of $z$.} That is the proportion of the data below $z$.

\vspace{0.3em}
{\setlength{\tabcolsep}{3.5pt}%
\begin{tabular}{|L{291pt}|L{198pt}|}\hline
\textbf{Step} & \textbf{Example: heights}\\ \hline
1. Draw the curve. Shade the part the question asks about. & Shorter than 68 inches: shade below 68.\\ \hline
2. Find each z-score: \; $z = \frac{x - \mu}{\sigma}$. Round to 2 decimal places. & $z = \frac{68 - 64}{3} = 1.33$\\ \hline
3. Look up each z-score in the table. & Row 1.3, column .03: .9082, so \textbf{90.82\%}\\ \hline
4. \textbf{Below:} use the table. \textbf{Above:} $1 -$ table. \newline \textbf{Between:} subtract the two table values. &
Taller than 68: $1 - .9082 = .0918$, \textbf{9.18\%} \newline 62 to 68: $z = \frac{62 - 64}{3} = -0.67$, table .2514. \newline $.9082 - .2514 = .6568$, \textbf{65.68\%}\\ \hline
\end{tabular}}

\sumhead{Normal Calculations in Reverse}
\begin{tabular}{|L{3.55in}|L{3.1in}|}\hline
\textbf{Step} & \textbf{Example: heights, 90th percentile}\\ \hline
1. Draw the curve. Shade the area to the \textbf{left}. & 90\% is to the left. The area is .9000.\\ \hline
2. Look \textbf{inside} the table for the area closest to it. & The closest area is .8997.\\ \hline
3. Read the z-score from the row and the column. & Row 1.2, column .08: \; $z = 1.28$\\ \hline
4. Find the value: \; $x = \mu + z\sigma$ & $64 + 1.28 \times 3 = 67.84$, about \textbf{67.8 inches}\\ \hline
\textbf{Bottom} 10\%: the area to the left is .10. & Closest area .1003, at $z = -1.28$. \newline $64 + (-1.28)(3) = 60.16$, about \textbf{60.2 inches}\\ \hline
\textbf{Top} 10\%: the area to the left is $1 - .10 = .90$. & Subtract from 1 \textbf{before} you look in the table.\\ \hline
\end{tabular}

\end{document}
